\documentclass[11pt]{article}
\usepackage{mathblatt}
% gesetzt von werkzeuge/setzer.py aus dem Lernweg (L3-A · L3-B · L3-C · L3-D · L3-T) und den Bankzeilen; Muster T6B.
% Präambel des Setzers (werkzeuge/setzer.py), wortgleich aus dem Hand-Satz T6B (bau/proben/2026-10-09/kathete-berechnen), dort aus K4W.
\makeatletter
\setlength{\mbpfbe}{0mm}% keine Punkte (Bauregel 6.4)
% eingerückt (Bauregel 4.7, 6.6): \gEin Einzug, \gSchrift eine Stufe kleiner
\newlength{\gEin}\setlength{\gEin}{0pt}
\let\gSchrift\relax
\renewcommand{\pfkreuzzeile}[1]{\par\noindent\hangindent3.2em\hangafter1\hspace*{1.6em}$\square$\ #1\par}
\newcommand{\gkreuz}[1]{$\square$\ #1\hspace{2em}}
% Bauregel 6.7: links, was man ansieht, rechts, was man tut; Spaltengrenze überall an derselben Stelle
\newsavebox{\gbox}\newlength{\gLinks}
\newcommand{\gzwei}[2]{\par\smallskip\noindent\sbox{\gbox}{#1}%
  \setlength{\gLinks}{\dimexpr0.5\textwidth-\gEin-\mbpfnr\relax}%
  \ifdim\wd\gbox>\dimexpr\gLinks-4mm\relax
    \usebox{\gbox}\par\smallskip\noindent\begin{minipage}{\linewidth}#2\end{minipage}%
  \else
    \begin{minipage}[t]{\gLinks}\vspace{0pt}\usebox{\gbox}\end{minipage}%
    \begin{minipage}[t]{\dimexpr\linewidth-\gLinks\relax}\vspace{0pt}\raggedright #2\end{minipage}%
  \fi\par}
\RenewDocumentEnvironment{pfaufg}{m m m}{%
  \begin{lrbox}{\mbaufgabebox}\begin{minipage}[t]{\linewidth}%
  \gSchrift\noindent\hspace*{\gEin}\mbpfrandmarke{#1}%
  \begin{minipage}[t]{\mbpfnr}\strut\textbf{#2}\end{minipage}%
  \begin{minipage}[t]{\dimexpr\linewidth-\gEin-\mbpfnr\relax}\raggedright\strut\ignorespaces}%
 {\end{minipage}%
  \end{minipage}\end{lrbox}%
  \setlength{\mbaufgabehoehe}{\dimexpr\ht\mbaufgabebox+\dp\mbaufgabebox\relax}%
  \par\addvspace{7pt}\Needspace*{\mbaufgabehoehe}%
  \noindent\usebox{\mbaufgabebox}\par\addvspace{7pt}}
\makeatother
% Rechenraum als Karo (Bauregel 6.9): n Rechenschritte = n mal 10 mm
\newcommand{\karo}[1]{\par\smallskip\noindent\begin{tikzpicture}
  \clip (0,0) rectangle ({\linewidth-1mm},{#1*10mm});
  \draw[black!16,line width=0.3pt,step=5mm] (0,0) grid ({\linewidth},{#1*10mm});
  \end{tikzpicture}\par}
\newcommand{\kk}[1]{$\square$\,#1\hspace{0.9em}}
\newcommand{\linie}[1]{\,{\color{black!45}\rule[-1pt]{#1}{0.4pt}}\,}
% Zeile am Ende jedes Durchgangs: Originale klein grau, darüber; dann Vorgänger/Nachfolger (Bauregel 3.4, 6.5)
\newcommand{\blattende}[2]{\par\vfill\noindent\begin{minipage}{\linewidth}{\scriptsize\color{mbgrau}Originale: #1\par}\smallskip
{\footnotesize #2\par}\end{minipage}}
\newcommand{\pkt}{\textperiodcentered{}}

\newcommand{\kennung}{FS9}
\newcommand{\grau}[1]{{\color{mbgrau}#1}}

\begin{document}
\pfheftstil
\fancyfoot[C]{\parbox[t]{\textwidth}{\mbfhbox\par\vspace{1.5mm}{\footnotesize\color{mbgrau}{\scriptsize \kennung}\hfill\thepage}}}
\pfheftkopf{Brüche vergleichen}{Klasse 6}
% L3-A: Gleicher Nenner oder gleicher Zähler
\begin{pfaufg}{}{1.}{}
\gzwei{\begin{tikzpicture}[font=\small]\fill[black!25] (0,0.00) rectangle (1.575,0.50);\draw[line width=0.4pt] (0.525,0.00) -- (0.525,0.50);\draw[line width=0.4pt] (1.050,0.00) -- (1.050,0.50);\draw[line width=0.4pt] (1.575,0.00) -- (1.575,0.50);\draw[line width=0.4pt] (2.100,0.00) -- (2.100,0.50);\draw[line width=0.4pt] (2.625,0.00) -- (2.625,0.50);\draw[line width=0.4pt] (3.150,0.00) -- (3.150,0.50);\draw[line width=0.4pt] (3.675,0.00) -- (3.675,0.50);\draw[line width=0.8pt] (0,0.00) rectangle (4.2,0.50);\node[left] at (0,0.25) {$\dfrac{3}{8}$};\fill[black!25] (0,-0.95) rectangle (2.625,-0.45);\draw[line width=0.4pt] (0.525,-0.95) -- (0.525,-0.45);\draw[line width=0.4pt] (1.050,-0.95) -- (1.050,-0.45);\draw[line width=0.4pt] (1.575,-0.95) -- (1.575,-0.45);\draw[line width=0.4pt] (2.100,-0.95) -- (2.100,-0.45);\draw[line width=0.4pt] (2.625,-0.95) -- (2.625,-0.45);\draw[line width=0.4pt] (3.150,-0.95) -- (3.150,-0.45);\draw[line width=0.4pt] (3.675,-0.95) -- (3.675,-0.45);\draw[line width=0.8pt] (0,-0.95) rectangle (4.2,-0.45);\node[left] at (0,-0.70) {$\dfrac{5}{8}$};\fill[black!25] (0,-1.90) rectangle (2.520,-1.40);\draw[line width=0.4pt] (0.840,-1.90) -- (0.840,-1.40);\draw[line width=0.4pt] (1.680,-1.90) -- (1.680,-1.40);\draw[line width=0.4pt] (2.520,-1.90) -- (2.520,-1.40);\draw[line width=0.4pt] (3.360,-1.90) -- (3.360,-1.40);\draw[line width=0.8pt] (0,-1.90) rectangle (4.2,-1.40);\node[left] at (0,-1.65) {$\dfrac{3}{5}$};\end{tikzpicture}}{\grau{a)\ Vergleiche $\dfrac{3}{8}$ mit $\dfrac{5}{8}$ und $\dfrac{3}{8}$ mit $\dfrac{3}{5}$.\par\smallskip $\dfrac{3}{8}$ und $\dfrac{5}{8}$: gleicher Nenner, beides Achtel. 5 Achtel sind mehr als 3 Achtel: $\dfrac{3}{8} < \dfrac{5}{8}$.\par $\dfrac{3}{8}$ und $\dfrac{3}{5}$: gleicher Zähler. Fünftel sind größer als Achtel – das Ganze ist in weniger Teile geteilt.\par 3 große Teile sind mehr als 3 kleine: $\dfrac{3}{5} > \dfrac{3}{8}$.}}
\par\medskip
b)\ Setze $<$ oder $>$ ein. \par\vspace{3pt} a)~$\dfrac{4}{9} \;\square\; \dfrac{7}{9}$ \qquad b)~$\dfrac{11}{12} \;\square\; \dfrac{5}{12}$ \qquad c)~$\dfrac{6}{10} \;\square\; \dfrac{9}{10}$\karo{3}
\fusshilfe{\mbox{1:~b) a) $<$ \quad b) $>$ \quad c) $<$}}
\end{pfaufg}

% L3-A: Gleicher Nenner oder gleicher Zähler
\begin{pfaufg}{}{2.}{}
Setze $<$ oder $>$ ein. \par\vspace{3pt} a)~$\dfrac{2}{3} \;\square\; \dfrac{2}{7}$ \qquad b)~$\dfrac{5}{12} \;\square\; \dfrac{5}{6}$ \qquad c)~$\dfrac{7}{10} \;\square\; \dfrac{7}{8}$\karo{3}
\fusshilfe{\mbox{2:~a) $>$ \quad b) $<$ \quad c) $<$}}
\end{pfaufg}

% L3-A: Gleicher Nenner oder gleicher Zähler
\begin{pfaufg}{}{3.}{}
Ordne von klein nach groß: \ $\dfrac{1}{4}$, \ $\dfrac{1}{10}$, \ $\dfrac{1}{3}$, \ $\dfrac{1}{6}$\karo{3}
\fusshilfe{\mbox{3:~$\frac{1}{10} < \frac{1}{6} < \frac{1}{4} < \frac{1}{3}$}}
\end{pfaufg}

% L3-A: Gleicher Nenner oder gleicher Zähler
\begin{pfaufg}{}{4.}{}
Jan, Pia und Ole laufen jeder eine gleich lange Runde um den See. Jan hat schon drei Viertel geschafft, Pia drei Fünftel und Ole drei Achtel. Wer ist am weitesten? Wer ist am wenigsten weit?\karo{3}
\fusshilfe{\mbox{4:~Jan am weitesten, Ole am wenigsten weit}}
\end{pfaufg}

% L3-B: Gleichnamig machen, dann vergleichen
\begin{pfaufg}{}{5.}{}
\grau{a)\ Was ist größer: $\dfrac{5}{6}$ oder $\dfrac{7}{8}$?\par\smallskip Gemeinsamer Nenner: 24 steht in der Sechser- und in der Achterreihe.\par Erweitern: $\dfrac{5}{6} = \dfrac{5 \cdot 4}{6 \cdot 4} = \dfrac{20}{24}$ \quad $\dfrac{7}{8} = \dfrac{7 \cdot 3}{8 \cdot 3} = \dfrac{21}{24}$\par Zähler vergleichen: $20 < 21$, also $\dfrac{5}{6} < \dfrac{7}{8}$.}
\par\medskip
b)\ Mach gleichnamig und setze $<$ oder $>$ ein. \par\vspace{3pt} a)~$\dfrac{2}{3} \;\square\; \dfrac{5}{6}$ \qquad b)~$\dfrac{3}{4} \;\square\; \dfrac{5}{8}$ \qquad c)~$\dfrac{7}{10} \;\square\; \dfrac{3}{5}$\karo{3}
\fusshilfe{\mbox{5:~b) a) $<$ \quad b) $>$ \quad c) $>$}}
\end{pfaufg}

% L3-B: Gleichnamig machen, dann vergleichen
\begin{pfaufg}{}{6.}{}
Mach gleichnamig und setze $<$ oder $>$ ein. \par\vspace{3pt} a)~$\dfrac{2}{3} \;\square\; \dfrac{3}{4}$ \qquad b)~$\dfrac{4}{5} \;\square\; \dfrac{3}{4}$ \qquad c)~$\dfrac{5}{6} \;\square\; \dfrac{7}{9}$\karo{3}
\fusshilfe{\mbox{6:~a) $<$ \quad b) $>$ \quad c) $>$}}
\end{pfaufg}

% L3-B: Gleichnamig machen, dann vergleichen
\begin{pfaufg}{}{7.}{}
Gleicher Nenner, gleicher Zähler oder keins von beiden? Entscheide erst, dann setze $<$ oder $>$ ein. \par\vspace{3pt} a)~$\dfrac{3}{4} \;\square\; \dfrac{9}{10}$ \qquad b)~$\dfrac{5}{7} \;\square\; \dfrac{5}{9}$ \qquad c)~$\dfrac{2}{5} \;\square\; \dfrac{3}{7}$\karo{3}
\fusshilfe{\mbox{7:~a) $<$ \quad b) $>$ \quad c) $<$}}
\end{pfaufg}

% L3-B: Gleichnamig machen, dann vergleichen
\begin{pfaufg}{}{8.}{}
Der Ausflug der 6b findet statt, wenn mindestens drei Viertel der Klasse zustimmen. 17 von 24 Kindern stimmen zu. Findet der Ausflug statt? Wenn nicht: Wie viele Stimmen fehlen?\karo{3}
\fusshilfe{\mbox{8:~Nein, es fehlt 1 Stimme.}}
\end{pfaufg}

% L3-C: Am Zahlenstrahl: ablesen, eintragen, ordnen
\begin{pfaufg}{}{9.}{}
\gzwei{\begin{tikzpicture}\draw[->,line width=0.7pt] (0,0) -- (5.75,0);\draw[line width=0.7pt] (0.000,-0.17) -- (0.000,0.17);\node[below] at (0.000,-0.17) {0};\draw[line width=0.4pt] (0.450,-0.1) -- (0.450,0.1);\draw[line width=0.4pt] (0.900,-0.1) -- (0.900,0.1);\draw[line width=0.4pt] (1.350,-0.1) -- (1.350,0.1);\draw[line width=0.4pt] (1.800,-0.1) -- (1.800,0.1);\draw[line width=0.4pt] (2.250,-0.1) -- (2.250,0.1);\draw[line width=0.7pt] (2.700,-0.17) -- (2.700,0.17);\node[below] at (2.700,-0.17) {1};\draw[line width=0.4pt] (3.150,-0.1) -- (3.150,0.1);\draw[line width=0.4pt] (3.600,-0.1) -- (3.600,0.1);\draw[line width=0.4pt] (4.050,-0.1) -- (4.050,0.1);\draw[line width=0.4pt] (4.500,-0.1) -- (4.500,0.1);\draw[line width=0.4pt] (4.950,-0.1) -- (4.950,0.1);\draw[line width=0.7pt] (5.400,-0.17) -- (5.400,0.17);\node[below] at (5.400,-0.17) {2};\fill (2.250,0) circle (1.7pt);\node[above] at (2.250,0.12) {$A$};\fill (0.900,0) circle (1.7pt);\node[above] at (0.900,0.12) {$\dfrac{1}{3}$};\fill (4.050,0) circle (1.7pt);\node[above] at (4.050,0.12) {$\dfrac{3}{2}$};\end{tikzpicture}}{\grau{a)\ Lies $A$ ab. Trage $\dfrac{1}{3}$ und $\dfrac{3}{2}$ ein und ordne alle drei.\par\smallskip Von 0 bis 1 sind es 6 gleiche Schritte: Ein Schritt ist $\dfrac{1}{6}$.\par $A$ liegt 5 Schritte hinter 0: $A = \dfrac{5}{6}$.\par $\dfrac{1}{3} = \dfrac{2}{6}$: 2 Schritte. \quad $\dfrac{3}{2} = \dfrac{9}{6}$: 9 Schritte, also hinter der 1.\par Von links nach rechts: $\dfrac{1}{3} < \dfrac{5}{6} < \dfrac{3}{2}$.}}
\par\medskip
\gzwei{\begin{tikzpicture}\draw[->,line width=0.7pt] (0,0) -- (5.75,0);\draw[line width=0.7pt] (0.000,-0.17) -- (0.000,0.17);\node[below] at (0.000,-0.17) {0};\draw[line width=0.4pt] (0.540,-0.1) -- (0.540,0.1);\draw[line width=0.4pt] (1.080,-0.1) -- (1.080,0.1);\draw[line width=0.4pt] (1.620,-0.1) -- (1.620,0.1);\draw[line width=0.4pt] (2.160,-0.1) -- (2.160,0.1);\draw[line width=0.7pt] (2.700,-0.17) -- (2.700,0.17);\node[below] at (2.700,-0.17) {1};\draw[line width=0.4pt] (3.240,-0.1) -- (3.240,0.1);\draw[line width=0.4pt] (3.780,-0.1) -- (3.780,0.1);\draw[line width=0.4pt] (4.320,-0.1) -- (4.320,0.1);\draw[line width=0.4pt] (4.860,-0.1) -- (4.860,0.1);\draw[line width=0.7pt] (5.400,-0.17) -- (5.400,0.17);\node[below] at (5.400,-0.17) {2};\fill (1.080,0) circle (1.7pt);\node[above] at (1.080,0.12) {$A$};\fill (2.160,0) circle (1.7pt);\node[above] at (2.160,0.12) {$B$};\fill (3.780,0) circle (1.7pt);\node[above] at (3.780,0.12) {$C$};\end{tikzpicture}}{%
b)\ Welche Brüche gehören zu $A$, $B$ und $C$?\karo{3}}
\fusshilfe{\mbox{9:~b) $A = \frac{2}{5}$, \ $B = \frac{4}{5}$, \ $C = \frac{7}{5}$}}
\end{pfaufg}

% L3-C: Am Zahlenstrahl: ablesen, eintragen, ordnen
\begin{pfaufg}{}{10.}{}
\gzwei{\begin{tikzpicture}\draw[->,line width=0.7pt] (0,0) -- (5.75,0);\draw[line width=0.7pt] (0.000,-0.17) -- (0.000,0.17);\node[below] at (0.000,-0.17) {0};\draw[line width=0.4pt] (0.900,-0.1) -- (0.900,0.1);\draw[line width=0.4pt] (1.800,-0.1) -- (1.800,0.1);\draw[line width=0.4pt] (2.700,-0.1) -- (2.700,0.1);\draw[line width=0.4pt] (3.600,-0.1) -- (3.600,0.1);\draw[line width=0.4pt] (4.500,-0.1) -- (4.500,0.1);\draw[line width=0.7pt] (5.400,-0.17) -- (5.400,0.17);\node[below] at (5.400,-0.17) {1};\end{tikzpicture}}{%
Der Zahlenstrahl ist in Sechstel geteilt. Trage $\dfrac{1}{3}$, $\dfrac{1}{2}$ und $\dfrac{5}{6}$ ein.\karo{3}}
\fusshilfe{\mbox{10:~$\frac{2}{6}$, \ $\frac{3}{6}$, \ $\frac{5}{6}$}}
\end{pfaufg}

% L3-C: Am Zahlenstrahl: ablesen, eintragen, ordnen
\begin{pfaufg}{}{11.}{}
Kleiner als $\dfrac{1}{2}$, zwischen $\dfrac{1}{2}$ und 1 oder größer als 1? Sortiere in drei Gruppen. \par\vspace{3pt} $\dfrac{3}{7}$ \qquad $\dfrac{5}{9}$ \qquad $\dfrac{8}{7}$ \qquad $\dfrac{2}{5}$ \qquad $\dfrac{6}{11}$ \qquad $\dfrac{11}{10}$\karo{3}
\fusshilfe{\mbox{11:~$\frac{3}{7}, \frac{2}{5}$ \,|\, $\frac{5}{9}, \frac{6}{11}$ \,|\, $\frac{8}{7}, \frac{11}{10}$}}
\end{pfaufg}

% L3-C: Am Zahlenstrahl: ablesen, eintragen, ordnen
\begin{pfaufg}{}{12.}{}
Ordne von klein nach groß: \ $\dfrac{5}{6}$, \ $\dfrac{2}{5}$, \ $\dfrac{7}{4}$, \ $\dfrac{3}{8}$\karo{3}
\fusshilfe{\mbox{12:~$\frac{3}{8} < \frac{2}{5} < \frac{5}{6} < \frac{7}{4}$}}
\end{pfaufg}

% L3-C: Am Zahlenstrahl: ablesen, eintragen, ordnen
\begin{pfaufg}{}{13.}{}
Vier Regentonnen sind verschieden hoch. \par $A$: 50\,cm hoch, 30\,cm Wasser \qquad $B$: 72\,cm hoch, 45\,cm Wasser \par $C$: 80\,cm hoch, 36\,cm Wasser \qquad $D$: 80\,cm hoch, 70\,cm Wasser \par Ordne die Tonnen danach, wie voll sie sind. Welche ist am vollsten?\karo{3}
\fusshilfe{\mbox{13:~$C < A < B < D$; $D$ am vollsten}}
\end{pfaufg}

% L3-D: Einen Bruch dazwischen finden
\begin{pfaufg}{}{14.}{}
\grau{a)\ Gib einen Bruch an, der zwischen $\dfrac{3}{5}$ und $\dfrac{2}{3}$ liegt.\par\smallskip Gleichnamig: $\dfrac{3}{5} = \dfrac{9}{15}$, \ $\dfrac{2}{3} = \dfrac{10}{15}$. Zwischen 9 und 10 liegt kein Zähler.\par Beide noch einmal mit 2 erweitern: $\dfrac{18}{30}$ und $\dfrac{20}{30}$.\par Dazwischen liegt $\dfrac{19}{30}$: \ $\dfrac{3}{5} < \dfrac{19}{30} < \dfrac{2}{3}$.}
\par\medskip
b)\ Gib zwei Brüche an, die dazwischen liegen. \par\vspace{3pt} a)~$\dfrac{2}{9}$ und $\dfrac{6}{9}$ \qquad b)~$\dfrac{5}{11}$ und $\dfrac{9}{11}$\karo{3}
\fusshilfe{\mbox{14:~b) a) z.\,B}}
\end{pfaufg}

% L3-D: Einen Bruch dazwischen finden
\begin{pfaufg}{}{15.}{}
Gib einen Bruch an, der dazwischen liegt. \par\vspace{3pt} a)~$\dfrac{3}{8}$ und $\dfrac{4}{8}$ \qquad b)~$\dfrac{1}{5}$ und $\dfrac{2}{5}$\karo{3}
\fusshilfe{\mbox{15:~a) $\frac{7}{16}$ \quad b) $\frac{3}{10}$}}
\end{pfaufg}

% L3-D: Einen Bruch dazwischen finden
\begin{pfaufg}{}{16.}{}
Gib einen Bruch an, der dazwischen liegt. \par\vspace{3pt} a)~$\dfrac{1}{3}$ und $\dfrac{3}{4}$ \qquad b)~$\dfrac{2}{5}$ und $\dfrac{1}{2}$\karo{3}
\fusshilfe{\mbox{16:~a) z.\,B}}
\end{pfaufg}

% L3-D: Einen Bruch dazwischen finden
\begin{pfaufg}{}{17.}{}
Welcher Bruch liegt zwischen $\dfrac{1}{4}$ und $\dfrac{2}{3}$? Kreuze an.\par\medskip\noindent\hspace*{7mm}\mbox{$\square$\ $\dfrac{1}{5}$}\hspace{10mm}\mbox{$\square$\ $\dfrac{3}{4}$}\hspace{10mm}\mbox{$\square$\ $\dfrac{5}{12}$}\hspace{10mm}\mbox{$\square$\ $\dfrac{7}{10}$}\rule[-3mm]{0pt}{8mm}\par\smallskip 
\fusshilfe{\mbox{17:~$\frac{5}{12}$}}
\end{pfaufg}

% L3-D: Einen Bruch dazwischen finden
\begin{pfaufg}{}{18.}{}
Mia braucht für eine Schleife ein Stück Band. Es soll länger als $\dfrac{1}{5}$\,m und kürzer als $\dfrac{1}{4}$\,m sein. Gib eine passende Länge als Bruch von 1\,m an. Wie viele Zentimeter sind das?\karo{3}
\fusshilfe{\mbox{18:~z.\,B. $\frac{9}{40}$\,m $=$ 22{,}5\,cm}}
\end{pfaufg}

% L3-T: Probetest
\begin{pfaufg}{}{19.}{}
Setze $<$ oder $>$ ein. \par\vspace{3pt} a)~$\dfrac{5}{9} \;\square\; \dfrac{4}{9}$ \qquad b)~$\dfrac{2}{7} \;\square\; \dfrac{2}{5}$ \qquad c)~$\dfrac{3}{4} \;\square\; \dfrac{7}{10}$\karo{3}
\fusshilfe{\mbox{19:~a) $>$ \quad b) $<$ \quad c) $>$}}
\end{pfaufg}

% L3-T: Probetest
\begin{pfaufg}{}{20.}{}
\gzwei{\begin{tikzpicture}\draw[->,line width=0.7pt] (0,0) -- (5.75,0);\draw[line width=0.7pt] (0.000,-0.17) -- (0.000,0.17);\node[below] at (0.000,-0.17) {0};\draw[line width=0.4pt] (0.675,-0.1) -- (0.675,0.1);\draw[line width=0.4pt] (1.350,-0.1) -- (1.350,0.1);\draw[line width=0.4pt] (2.025,-0.1) -- (2.025,0.1);\draw[line width=0.7pt] (2.700,-0.17) -- (2.700,0.17);\node[below] at (2.700,-0.17) {1};\draw[line width=0.4pt] (3.375,-0.1) -- (3.375,0.1);\draw[line width=0.4pt] (4.050,-0.1) -- (4.050,0.1);\draw[line width=0.4pt] (4.725,-0.1) -- (4.725,0.1);\draw[line width=0.7pt] (5.400,-0.17) -- (5.400,0.17);\node[below] at (5.400,-0.17) {2};\fill (2.025,0) circle (1.7pt);\node[above] at (2.025,0.12) {$A$};\fill (3.375,0) circle (1.7pt);\node[above] at (3.375,0.12) {$B$};\end{tikzpicture}}{%
Welche Brüche gehören zu $A$ und $B$? Trage $\dfrac{3}{2}$ ein.\karo{3}}
\fusshilfe{\mbox{20:~$A = \frac{3}{4}$, \ $B = \frac{5}{4}$, \ $\frac{3}{2} = \frac{6}{4}$}}
\end{pfaufg}

% L3-T: Probetest
\begin{pfaufg}{}{21.}{}
Ordne von klein nach groß: \ $\dfrac{4}{7}$, \ $\dfrac{1}{3}$, \ $\dfrac{5}{4}$, \ $\dfrac{1}{2}$\karo{3}
\fusshilfe{\mbox{21:~$\frac{1}{3} < \frac{1}{2} < \frac{4}{7} < \frac{5}{4}$}}
\end{pfaufg}

% L3-T: Probetest
\begin{pfaufg}{}{22.}{}
Gib eine Zahl an, die zwischen $\dfrac{3}{5}$ und $\dfrac{5}{6}$ liegt.\karo{3}
\fusshilfe{\mbox{22:~z.\,B. $\frac{2}{3}$}}
\end{pfaufg}

% L3-T: Probetest
\begin{pfaufg}{}{23.}{}
Welche Aussage ist wahr? Kreuze an.\par\medskip\noindent\hspace*{7mm}\mbox{$\square$\ $\dfrac{3}{4} > \dfrac{9}{10}$}\hspace{10mm}\mbox{$\square$\ $\dfrac{5}{8} < \dfrac{1}{2}$}\hspace{10mm}\mbox{$\square$\ $\dfrac{7}{12} > \dfrac{1}{2}$}\hspace{10mm}\mbox{$\square$\ $\dfrac{2}{9} > \dfrac{2}{7}$}\rule[-3mm]{0pt}{8mm}\par\smallskip 
\fusshilfe{\mbox{23:~$\frac{7}{12} > \frac{1}{2}$}}
\end{pfaufg}

% L3-T: Probetest
\begin{pfaufg}{}{24.}{}
Der Förderverein belohnt die Klasse, in der ein größerer Anteil der Kinder beim Spendenlauf mitläuft. In der 6a laufen 21 von 28 Kindern mit, in der 6b 20 von 25 Kindern. Welche Klasse wird belohnt?\karo{3}
\fusshilfe{\mbox{24:~die 6b}}
\end{pfaufg}

\par\vfill\noindent{\footnotesize Vorher: Kürzen und Erweitern\quad\textbullet\quad Weiter: Dezimalzahlen\par}
\end{document}
